% systemd-deep.tex — how systemd thinks: units and their names (a dash is a slash), the two
% dependency axes, Type= as "when is a start done", Restart= and the start limit, stopping by cgroup,
% PID 1 that may not die, socket activation (the kernel buffer replaces ordering), sd_notify and the
% fd store, timers vs cron, the slice / scope / service tree, the sandbox as a namespace view,
% DynamicUser's leased UIDs, credentials, the user manager, journald's trusted fields and sealing.
% Pictures: one service on both axes · the service state machine · the Type= timeline · the Restart=
% matrix · the socket buffer · the cgroup tree · what a sandboxed service sees · a release timeline.
% Sources (checked 2026-10-10):
%   man7.org (pages carrying systemd 262~devel): systemd.unit(5) (unit types; load path, first match
%     wins; drop-ins merged alphanumerically, dash-truncated foo-.service.d; masked = empty or
%     /dev/null, "cannot be activated"; string escaping: "/" -> "-", other non-alnum (not : _ .)
%     -> C-style "\x2d", root "/" -> "-"; Wants / Requires ("fails only with After="; stop / restart
%     propagate) / Requisite / BindsTo / PartOf / Upholds 249 / Conflicts; OnFailure, OnSuccess 249;
%     ordering orthogonal; stop = inverse; Condition skips, Assert fails; specifiers %i %n %u %h)
%   · systemd.service(5) (simple "will report success even if the service's binary cannot be
%     invoked"; "recommended to use Type=exec for long-running services"; exec = after execve();
%     notify-reload 253; Restart= table; on-failure "the recommended choice for long-running
%     services"; RestartSec 100 ms; RestartSteps / RestartMaxDelaySec 254; RestartRandomizedDelaySec
%     262; RestartMode=direct; "|" prefix = the user's default shell; WatchdogSec -> SIGABRT;
%     SuccessExitStatus 0 + SIGHUP SIGINT SIGTERM SIGPIPE)
%   · systemd.exec(5) (DynamicUser 232: "allocated from the UID/GID range 61184...65519"; implies
%     RemoveIPC, PrivateTmp, NoNewPrivileges, RestrictSUIDSGID, ProtectSystem=strict,
%     ProtectHome=read-only; StateDirectory under /var/lib/private, mode 0700, symlinked; "Sets up a
%     new file system namespace"; ProtectSystem=strict "entire file system hierarchy is mounted
%     read-only, except for /dev/, /proc/ and /sys/"; ProtectHome "/home/, /root, and /run/user are
%     made inaccessible and empty"; PrivateNetwork "new network namespace … only the loopback";
%     PrivateDevices; Environment "exposed to unprivileged clients via D-Bus IPC"; LoadCredential 247
%     "read-only location … non-swappable memory … only accessible to the user associated with the
%     unit"; LoadCredentialEncrypted 250)
%   · systemd(1) (systemd.crash_action= "freeze", "reboot", "poweroff", default freeze: "the system
%     will hang indefinitely when the system manager (PID 1) crashes"; crash_shell; SIGTERM to PID 1:
%     "serializes its state, reexecutes itself and deserializes the saved state again … equivalent to
%     systemctl daemon-reexec")
%   · systemd-system.conf(5) (DefaultTimeoutStopSec 90 s; StartLimit 10 s / 5; DefaultTasksMax 15 %)
%   · systemd.kill(5) (KillMode control-group default, mixed, process / none discouraged)
%   · systemd.slice(5) ("dash-separated series of names … path to the slice from the root slice";
%     root -.slice; foo-bar.slice inside foo.slice) · systemd-escape(1)
%   · systemd.timer(5) (AccuracySec 1 min: "host-specific, randomized, but stable position that is
%     synchronized between all local timer units", to suppress wake-ups; RandomizedDelaySec;
%     Persistent only for OnCalendar, last trigger stored on disk; active unit "not restarted, but
%     simply left running")
%   · systemd.socket(5) (Accept= default no, "preferable"; template per connection; MaxConnections
%     64) · sd_listen_fds(3) (SD_LISTEN_FDS_START 3)
%   · sd_notify(3) ("a single datagram" to $NOTIFY_SOCKET, "newline-separated list of variable
%     assignments"; FDSTORE=1 held "and … handed back … at the next start or restart"; "possible to
%     reimplement … without external dependencies", stable)
%   · systemd.special(7) (network.target "only very weakly defined"; -.slice, system.slice,
%     user.slice, machine.slice, init.scope) · user@.service(5) (TasksMax=33%) · pam_systemd(8) ·
%     loginctl(1) (enable-linger) · systemctl(1) (edit → drop-in; enable makes [Install] symlinks, does
%     not start; soft-reboot 254) · systemd-run(1) · systemd-analyze(1) (blame "might be
%     misleading"; critical-chain; security 0.0-10.0) · systemd-sysext(8) (248; confext 254) · systemd.resource-control(5) (CPUWeight
%     1-10000, 100, idle; CPUQuota; TasksMax; Delegate=)
%   · journald.conf(5) (Storage= "Defaults to 'persistent' … determined at compilation time"; 10 % /
%     15 % capped at 4G; 10000 messages in 30 s; Seal= on by default when a key exists, FSS based on
%     "Seekable Sequential Key Generators", doi:10.1007/978-3-642-40203-6_7)
%   · journalctl(1) (--setup-keys: sealing key "shall remain on the host", verification key "stored
%     externally"; --interval= "Defaults to 15min", shorter "shorten the time range of undetectable
%     journal alterations"; --verify)
%   · systemd.journal-fields(7) ("Fields prefixed with an underscore are trusted fields … cannot be
%     altered by client code"; _BOOT_ID 128-bit)
%   eprint.iacr.org/2013/397 (Practical Secure Logging: Seekable Sequential Key Generators, ESORICS
%     2013; "our SSKG implementation has become part of the logging service of the systemd")
%   github.com/torvalds/linux kernel/exit.c (do_exit: is_global_init → panic("Attempted to kill
%     init! exitcode=…"))
%   github.com/systemd/systemd src/core/manager.c ("%sA %s job is running for %s (%s / %s)")
%   0pointer.de/blog/projects/systemd.html ("Rethinking PID 1", 2010-04-30: syslog clients' messages
%     "added to the /dev/log socket buffer … will not have to wait"; launchd "listening of the sockets
%     is pulled out of all daemons"; "Prior to launchd the venerable inetd worked much like this")
%   0pointer.net/blog/dynamic-users-with-systemd.html (2017-10-06: distributions allocate "below the
%     60000 range"; 65534 nobody, 65535 invalid; /var/lib/private root:root 0700 boundary)
%   brand.systemd.io (Spelling: lower-case d of a system daemon; "System Five Hundred since D is the
%     roman numeral for 500 (this also clarifies the relation to System V, right?)")
%   systemd.io/CGROUP_DELEGATION ("each cgroup only has a single writer"; slices = "trunk and
%     branches", scopes and services = leaves; a delegated service moves its main process into a
%     sub-cgroup)
%   0pointer.de/blog/projects/the-biggest-myths.html (2013-01-26)
%   api.github.com/repos/systemd/systemd/releases (v247 2020-11-26 · v250 2021-12-23 · v254
%     2023-07-28 · v256 2024-06-11 (run0; refuses cgroup v1) · v258 2025-09-17 (cgroup v1, telinit,
%     runlevel removed; "|" prefix) · v259 2025-12-17 (journal persistent) · v260 2026-03-17 (SysV
%     scripts, rc-local removed) · v262 2026-09-22)
%   github.com/systemd/systemd/issues/39338 ("deleted to break ordering cycle": often not the culprit)
% Not repeated: boot chain, target order, rescue / emergency, UKI, systemd-analyze time
% (linux-boot-process); Memory*= knobs, OOMPolicy=, systemd-oomd (linux-memory-oom); cgroup v2
% files and the v1 sunset of the container stack (namespaces-cgroups-containers); capability sets,
% no_new_privs, seccomp internals, run0 + polkit (linux-security-permissions).
% @labels: area=linux kind=tooling platform=backend topic=platform-apis,tooling,debugging discussion=321
% @tags: systemd, pid-1, unit-files, dependencies, ordering, service-type, restart-policy, drop-ins, socket-activation, timers, journald, cgroup-v2, slices, sandboxing, no-new-privs, seccomp, credentials, systemd-analyze
\documentclass[8pt]{extarticle}
\usepackage{printup-sheet}
\usepackage{array}

% Linked text: the website links it, paper prints the text only (paper holds no links).
\newcommand\link[2]{\pdfonly{#1}\htmlonly{\href{#2}{#1}}}

% unit-file listing — literate kills the mono font's ligatures
\lstdefinelanguage{UnitSheet}{
  sensitive=true, morecomment=[l]{\#},
  literate={--}{{\hbox{-}\hbox{-}}}2 {->}{{\hbox{-}\hbox{>}}}2 {==}{{\hbox{=}\hbox{=}}}2
           {!=}{{\hbox{!}\hbox{=}}}2 {>=}{{\hbox{>}\hbox{=}}}2 {<=}{{\hbox{<}\hbox{=}}}2}
\lstset{basicstyle=\ttfamily\scriptsize, aboveskip=2pt, belowskip=2pt}

\newcommand\ct[1]{\texttt{#1}}
\newcommand\dd{\hbox{-}\hbox{-}}          % "--" without the ligature
\newcommand\hd[1]{\par\vspace{2pt}\noindent{\small\bfseries\color{sheetBlue}#1}\par\vspace{1pt}}
\newcolumntype{L}[1]{>{\raggedright\arraybackslash}p{#1}}

\tikzset{
  u/.style={box, font=\scriptsize\ttfamily, inner sep=2pt, minimum height=5mm},
  st/.style={box, font=\scriptsize, inner sep=2pt, minimum height=5mm, minimum width=15mm},
  lbl/.style={font=\tiny, text=black!75, inner sep=1pt, align=center},
  dep/.style={->, thick, draw=sheetBlue},
  ord/.style={->, thick, dashed, draw=sheetOrange},
  bad/.style={->, thick, draw=sheetRed},
  cg/.style={draw=#1, fill=#1!8, rounded corners=1.5pt, font=\tiny\ttfamily, inner sep=1.2pt,
             anchor=west},
  why/.style={font=\tiny, anchor=west, inner sep=1pt, text=black!85},
  dir/.style={font=\tiny\ttfamily, anchor=west, inner sep=1pt, text=sheetGrey},
}

\begin{document}

\sheettitle{systemd — units, dependencies, services, sandbox, journal}{linux · folio}

\oneliner{systemd (PID~1) turns \textbf{unit files} into a \textbf{transaction of jobs} on two
independent axes — \textbf{requirement} (\ct{Wants}, \ct{Requires}: \emph{what} starts) and
\textbf{ordering} (\ct{After}, \ct{Before}: \emph{when}) — and gives each service its own
\textbf{cgroup} to limit, log and kill \emph{all} its processes.}

\vspace{1pt}
\noindent\begin{tikzpicture}[sheet]
  % ── dependency vs ordering graph ──
  \node[font=\bfseries\scriptsize, anchor=west, text=sheetBlue] at (-0.1,3.15) {One service, two axes: requirement (\textcolor{sheetBlue}{solid}) $\neq$ ordering (\textcolor{sheetOrange}{dashed})};
  \node[u, fill=sheetOrange!12, draw=sheetOrange] (app) at (3.9,1.45) {app.service};
  \node[u] (tgt) at (3.9,2.6) {multi-user.target};
  \node[u] (pg)  at (0.85,1.45) {postgresql.service};
  \node[u] (rd)  at (0.85,0.4) {redis.service};
  \node[u] (mnt) at (6.95,1.45) {srv-data.mount};
  \node[u] (sock) at (6.95,0.4) {app.socket};
  \node[u, fill=black!4, draw=sheetGrey] (log) at (3.9,0.4) {app-log.service};
  \draw[dep] (tgt) -- node[lbl, right]{\ct{Wants} (\ct{enable})} (app);
  \draw[dep] ([yshift=2pt]app.west) -- node[lbl, above]{\ct{Requires}} ([yshift=2pt]pg.east);
  \draw[ord] ([yshift=-2pt]app.west) -- node[lbl, below, text=sheetOrange!80!black]{\ct{After}} ([yshift=-2pt]pg.east);
  \draw[dep] (app.south west) -- node[lbl, above, sloped, pos=0.62]{\ct{Wants}} (rd.north east);
  \draw[dep] ([yshift=2pt]app.east) -- node[lbl, above]{\ct{BindsTo}} ([yshift=2pt]mnt.west);
  \draw[ord] ([yshift=-2pt]app.east) -- node[lbl, below, text=sheetOrange!80!black]{\ct{After}} ([yshift=-2pt]mnt.west);
  \draw[dep, draw=sheetGreen!60!black] (sock.north west) -- node[lbl, above, sloped, pos=0.38]{activates} (app.south east);
  \draw[dep, draw=sheetGrey] (log) -- node[lbl, right, pos=0.45]{\ct{PartOf}} (app);
  \node[lbl, anchor=north west, align=left, text=sheetRed] at (-0.1,2.9) {redis: \ct{Wants}, no \ct{After} $\to$ starts in\\\textbf{parallel}; its failure is ignored};
  \node[lbl, anchor=north east, align=right, text=sheetGreen!40!black] at (7.85,2.9) {pg start fails $\to$ app's job \textbf{fails}\\mount goes away $\to$ app \textbf{stops}};
  \node[lbl, anchor=north, text=sheetGrey] at (3.9,0.08) {stop / restart app $\to$ app-log follows};
  \draw[sheetGrey!50] (8.05,3.25) -- (8.05,-0.15);
  % ── service state machine ──
  \node[font=\bfseries\scriptsize, anchor=west, text=sheetBlue] at (8.2,3.15) {Service states (``Active:'' in \ct{systemctl status})};
  \node[st, fill=black!5, draw=sheetGrey] (in) at (9.05,2.45) {inactive};
  \node[st, text width=21mm] (act) at (11.6,2.45) {\textbf{activating}\\{\tiny ExecCondition, \dots Pre, ExecStart}};
  \node[st, fill=sheetGreen!12, draw=sheetGreen!60!black, text width=15mm] (on) at (14.95,2.45) {\textbf{active}\\{\tiny running · exited}};
  \node[st, text width=17mm] (deact) at (14.95,1.15) {deactivating\\{\tiny ExecStop, SIGTERM $\to$ inactive}};
  \node[st, fill=sheetRed!10, draw=sheetRed, text width=15mm] (fail) at (11.6,1.15) {\textbf{failed}\\{\tiny \ct{OnFailure=} runs}};
  \node[st, fill=sheetOrange!10, draw=sheetOrange, text width=15mm] (ar) at (9.05,1.15) {activating\\{\tiny (auto-restart)}};
  \draw[flow] (in) -- node[lbl, above]{start} (act);
  \draw[flow] (act) -- node[lbl, above]{\ct{Type=}: ready} (on);
  \draw[flow] (on) -- node[lbl, right]{stop} (deact);
  \draw[flow] ([xshift=2mm]on.north west) to[out=110, in=70, looseness=1.1] node[lbl, above]{reload: \ct{reloading}} ([xshift=-2mm]on.north east);
  \draw[bad] (act) -- node[lbl, right, text=sheetRed]{fail} (fail);
  \draw[bad] (on.south west) -- node[lbl, below, sloped, text=sheetRed, pos=0.4]{crash} (fail.east);
  \draw[hot] (fail) -- node[lbl, above]{\ct{Restart=}} node[lbl, below]{matches} (ar);
  \draw[hot] (ar.north east) -- node[lbl, above, sloped]{\ct{RestartSec}} (act.south west);
  \node[lbl, anchor=north west, align=left, text=black!80] at (8.2,0.6) {\textbf{fail} = exit $\neq$ 0, unclean signal, timeout, watchdog; \ct{Restart=no} (default): stays \textbf{failed}.\\
    $>$\ct{StartLimitBurst} (5) starts in \ct{StartLimitIntervalSec} (10\,s), manual ones too $\to$\\
    \textcolor{sheetRed}{\ct{start-limit-hit}}: no more restarts until \ct{reset-failed}. \ct{RestartMode=direct} (254):\\
    skips \textbf{failed}, so \ct{OnFailure=} does not run and dependants never see it.};
\end{tikzpicture}

\begin{multicols}{2}
\raggedright\footnotesize\setstretch{1.0}\setlength{\parskip}{1.5pt plus 1pt}

\hd{Units — and a dash in the name is a slash}
Every \link{unit}{https://man7.org/linux/man-pages/man5/systemd.unit.5.html} is typed by its suffix
(\ct{.service}, \ct{.mount}, \ct{.slice} … eleven kinds) and its name is an address: \ct{/srv/data} $\to$ \ct{srv-data.mount}; the root \ct{/} is a lone dash,
hence \ct{-.mount} and the root cgroup \ct{-.slice}; \ct{system-app.slice} sits \emph{inside}
\ct{system.slice}; \ct{foo-bar.service} merges drop-in \ct{*.conf}s from
\ct{foo-bar.service.d/} \emph{and} \ct{foo-.service.d/}. A real dash must be
\link{escaped}{https://man7.org/linux/man-pages/man1/systemd-escape.1.html}:
\ct{/dev/disk/by-label/x} $\to$ \ct{dev-disk-by\textbackslash x2dlabel-x.device}.
Layers, first match wins: \textbf{\ct{/etc}} (admin) $>$ \ct{/run} (runtime; generators turn
\ct{fstab} into \ct{.mount}s) $>$ \textbf{\ct{/usr/lib}} (package: override, never edit). Symlinked
to \ct{/dev/null}: \textbf{masked} — unstartable, even as a dependency. \ct{\#} comments only at line
start (\ct{KEY=v \#x} keeps ``\ct{\#x}''); an empty \ct{KEY=} resets a list.

\hd{How tightly X holds Y — loose to tight}
{\scriptsize\setlength\tabcolsep{2pt}
\begin{tabular}{@{}>{\ttfamily}L{12mm}L{66mm}@{}}
\toprule
Wants= & pulls Y in; Y's failure is ignored — the recommended default \\
Requires= & + Y's failed start fails X, \emph{only if} X is also \ct{After=} Y; stopping Y stops X \\
Requisite= & never starts Y: Y must already be up, or X fails at once \\
BindsTo= & + Y going away for \emph{any} reason (crash, unplugged disk) stops X \\
PartOf= & only Y's stop / restart reaches X — one way, nothing else \\
Upholds= & (249) the reverse pull: while X is up, a dead Y is started again \\
Conflicts= & starting one stops the other \\
OnFailure= & listed units start when X fails · \ct{OnSuccess=} (249) \\
\bottomrule
\end{tabular}}
\begin{itemize}
  \item \textbf{Ordering is orthogonal}: no \ct{After=}/\ct{Before=} $\Rightarrow$ \emph{parallel};
        stops run in inverse order, before starts. Every service silently gets
        \ct{After=basic.target} + \ct{Conflicts=shutdown.target}: that is how shutdown stops it.
  \item \ct{network.target} is ``only very weakly defined'': \ct{Wants=} +
        \ct{After=network-online.target}. A \textbf{cycle} is broken by deleting a job, often an
        innocent one's. \ct{Condition*=} false skips; \ct{Assert*=} fails.
\end{itemize}

\hd{\link{\ct{Type=}}{https://man7.org/linux/man-pages/man5/systemd.service.5.html} — the moment a start counts as done}
\noindent\begin{tikzpicture}[sheet, x=1cm, y=1cm]
  \draw[->, thick, sheetGrey] (0,0) -- (7.9,0) node[lbl, anchor=west]{t};
  \foreach \x/\t/\e/\k in {0.4/simple/fork()/sheetRed, 1.55/exec/execve()/sheetGreen!50!black,
      2.85/forking/parent exits/sheetGrey, 4.2/notify/READY=1 sent/sheetBlue,
      5.55/dbus/bus name taken/sheetBlue, 6.95/oneshot/process exits/sheetOrange} {
    \fill[\k] (\x,0) circle (0.6mm);
    \node[font=\scriptsize\ttfamily\bfseries, text=\k, anchor=south, inner sep=1pt] at (\x,0.05) {\t};
    \node[font=\tiny, anchor=north, inner sep=1pt] at (\x,-0.05) {\e}; }
\end{tikzpicture}\par
At its dot \ct{systemctl start} returns and \ct{After=} units begin. \ct{simple} is ``started''
before the binary even loads — a mistyped path still reports success; \ct{exec} (the man page's
pick) catches that, \ct{notify} waits for real readiness.

\hd{\ct{Restart=} — which exits restart}
\noindent\begin{tikzpicture}[sheet, x=1cm, y=1cm]
  \foreach \c [count=\i from 0] in {no, always, on-success, on-failure, on-abnormal, on-abort, on-watchdog}
    \node[font=\tiny\ttfamily, rotate=24, anchor=south west, inner sep=0.5pt] at (2.05+0.64*\i,0.02) {\c};
  \foreach \r [count=\j from 0] in {clean exit code or signal, unclean exit code, unclean signal, timeout, watchdog} {
    \node[font=\tiny, anchor=east, inner sep=1pt] at (2.0,-0.16-0.22*\j) {\r};
    \draw[sheetGrey!40] (2.0,-0.05-0.22*\j) -- (6.5,-0.05-0.22*\j); }
  \draw[sheetGrey!40] (2.0,-1.15) -- (6.5,-1.15);
  \fill[sheetGreen!60!black, opacity=0.12] (2.0+0.64*3,-1.15) rectangle (2.64+0.64*3,-0.05);
  \foreach \i/\j in {1/0,2/0, 1/1,3/1, 1/2,3/2,4/2,5/2, 1/3,3/3,4/3, 1/4,3/4,4/4,6/4}
    \fill[sheetBlue] (2.32+0.64*\i,-0.16-0.22*\j) circle (0.06);
  \node[font=\tiny, anchor=north west, align=left, text=sheetGreen!40!black, inner sep=0pt] at (6.62,-0.12) {\ct{on-failure}: ``the\\recommended\\choice for long-\\running services''};
\end{tikzpicture}\par
\textbf{Clean} = exit 0, or SIGHUP, SIGINT, SIGTERM, SIGPIPE (not for oneshot). Never after
\ct{systemctl stop}. \ct{RestartSec=} 100\,ms; back-off (254), jitter (262).

\hd{Stopping — the cgroup is the process list}
\ct{ExecStop=} $\to$ SIGTERM to \emph{every} process in the cgroup (\ct{KillMode=control-group},
the default: a \ct{nohup cmd \&} or double-forked child dies too; \ct{process}, \ct{none} let them
escape) $\to$ \ct{TimeoutStopSec} 90\,s $\to$ SIGKILL — the 90\,s of the shutdown line
\ct{A stop job is running for … (12s / 1min 30s)}. \ct{ExecStopPost=} runs even after a crash, told
why in \ct{\$SERVICE\_RESULT}. No \ct{WATCHDOG=1} within \ct{WatchdogSec=} $\to$ \textbf{SIGABRT}.
\ct{ExecStart=} is \emph{not} a shell line — no pipes, no \ct{>}; prefix \ct{|} (258) asks for one.

\hd{PID~1, which is not allowed to die}
If PID~1 exits, the kernel panics (\ct{Attempted to kill init!}). So a crashed systemd \textbf{hangs
the machine on purpose}:
\link{\ct{systemd.crash\_action=freeze}}{https://man7.org/linux/man-pages/man1/systemd.1.html} is the
default (or \ct{reboot}, \ct{poweroff}, a crash shell). \ct{kill -TERM 1} does not stop it either — it
\textbf{serialises its state, re-executes itself} and reads the state back, like \ct{systemctl
daemon-reexec}: how an upgraded systemd takes over without a reboot. The name: a system
\textbf{d}aemon, lower case; the project's own joke calls it ``System Five Hundred'' — D is Roman 500,
which ``clarifies the relation to System V''.

\hd{systemctl, systemd-analyze — the traps}
\ct{enable} only makes the \ct{[Install]} symlinks (\ct{WantedBy=} $\to$ a \ct{.wants/} link); it
does \emph{not} start (\ct{\dd now} does) · \ct{edit} writes a drop-in; to replace the command, first
an empty \ct{ExecStart=} — two get the unit refused · \ct{systemd-analyze blame} counts waiting as
start time (``might be misleading''); \ct{critical-chain} shows who waited for whom.

\hd{User manager, transient units, images}
\ct{systemctl \dd user} talks to \ct{user@1000.service}; it ends with the \textbf{last session} unless
\ct{loginctl enable-linger} (then: from boot, past logout). \ct{systemd-run cmd}: a throwaway
(transient) service; \ct{run0} (256) is the same binary. \ct{sysext} (248) / \ct{confext} overlay
\ct{/usr}, \ct{/etc} from images; \ct{soft-reboot}: new userspace, same kernel.

\columnbreak

\hd{Sockets first — the kernel's buffer replaces ordering}
\noindent\begin{tikzpicture}[sheet, x=1cm, y=1cm]
  \node[u] (c1) at (0.55,0.2) {logger};
  \node[u] (c2) at (0.55,-0.3) {sshd};
  \node[box, minimum width=20mm, minimum height=5mm] (buf) at (3.35,-0.05) {};
  \foreach \i in {0,...,5} \fill[sheetBlue!40] (2.5+0.17*\i,-0.2) rectangle (2.62+0.17*\i,0.1);
  \node[lbl, anchor=south] at (buf.north) {\ct{/dev/log}: bound by PID~1 first};
  \node[u, fill=sheetOrange!12, draw=sheetOrange] (sys) at (6.6,-0.05) {syslog \textrm{starting…}};
  \draw[flow] (c1.east) -- ([yshift=1pt]buf.west);
  \draw[flow] (c2.east) -- ([yshift=-1pt]buf.west);
  \draw[flow, dashed] (buf.east) -- node[lbl, above]{reads when up} (sys.west);
\end{tikzpicture}\par
PID~1 binds every listening socket \emph{before} any daemon runs, then starts them all at once: a
client writing while syslog is still starting just queues in the socket buffer — no \ct{After=}
needed. An old idea: \ct{inetd} did it per connection, Apple's launchd for every daemon.
\ct{Accept=yes} is inetd mode (an \ct{app@.service} per connection, max 64); \ct{Accept=no}
(default) hands over the listener as fd \textbf{3}, still bound across restarts.

\hd{sd\_notify, the fd store, templates}
\link{\ct{sd\_notify}}{https://man7.org/linux/man-pages/man3/sd_notify.3.html} is one datagram of
\ct{KEY=VALUE} lines to the socket named in \ct{\$NOTIFY\_SOCKET} — a stable protocol, a few lines
without libsystemd: \ct{READY=1} · \ct{RELOADING=1} · \ct{STOPPING=1} · \ct{STATUS=} ·
\ct{WATCHDOG=1}. \ct{FDSTORE=1} parks open fds \emph{in PID~1}, handed back at the next start: a
restart (or soft-reboot) that keeps its connections. \ct{app@.service} + \ct{start app@8080}
$\Rightarrow$ \ct{\%i}=\ct{8080}; \textbf{\ct{\%u}, \ct{\%h} are the manager's user} — \ct{root}
in a system unit, not \ct{User=}.

\hd{Timers — cron, minus its surprises}
A \link{\ct{.timer}}{https://man7.org/linux/man-pages/man5/systemd.timer.5.html} starts its
same-named service. Still running when it fires $\Rightarrow$ \textbf{no second copy};
\ct{Persistent=true} runs an \ct{OnCalendar=} job missed while the box was off;
\ct{RandomizedDelaySec=} spreads a fleet's 3 a.m. stampede. \ct{AccuracySec=} (1\,min) is deliberate
slack: all timers fire at one host-specific, stable point of that window, so the CPU wakes once
for many.

\hd{The cgroup tree, resources, delegation}
\noindent\begin{tikzpicture}[sheet, x=1cm, y=1cm]
  \node[cg=sheetGrey] (r) at (0,0) {-.slice};
  \node[cg=sheetGrey] (i) at (0.3,-0.27) {init.scope \textrm{PID 1}};
  \node[cg=sheetBlue] (s) at (0.3,-0.54) {system.slice};
  \node[cg=sheetOrange] (a) at (0.6,-0.81) {app.service};
  \node[cg=sheetBlue] (sa) at (0.6,-1.08) {system-app.slice};
  \node[cg=sheetOrange] (a8) at (0.9,-1.35) {app@8080.service};
  \node[cg=sheetBlue] (us) at (3.75,0) {user.slice};
  \node[cg=sheetBlue] (u1) at (4.05,-0.27) {user-1000.slice \textrm{TasksMax 33\%}};
  \node[cg=sheetGreen!60!black] (se) at (4.35,-0.54) {session-3.scope \textrm{ssh}};
  \node[cg=sheetOrange] (um) at (4.35,-0.81) {user@1000.service \textrm{Delegate}};
  \node[cg=sheetBlue] (m) at (3.75,-1.35) {machine.slice \textrm{VMs, containers}};
  \foreach \p/\c in {r/i, r/s, s/a, s/sa, sa/a8, us/u1, u1/se, u1/um}
    \draw[sheetGrey] ([xshift=1.2mm]\p.south west) |- (\c.west);
  \draw[sheetGrey] (r.east) -- (us.west);
  \draw[sheetGrey] (3.5,0) |- (m.west);
  \node[lbl, anchor=west, align=left] at (2.35,-0.95) {\textcolor{sheetBlue}{slices}: inner\\nodes, no\\processes};
\end{tikzpicture}\par
Services and scopes are the leaves. \ct{CPUWeight=} 1–10\,000 (100, or \ct{idle}): a share under
contention · \ct{CPUQuota=150\%} = 1.5 CPUs · \ct{TasksMax=} (15\,\% of the pid limit), live
via \ct{set-property}. \textbf{\link{\ct{Delegate=yes}}{https://systemd.io/CGROUP_DELEGATION/}}: systemd keeps out
of the subtree (one writer per cgroup); the service moves itself into a sub-cgroup.

\hd{What a sandboxed service sees}
\noindent\begin{tikzpicture}[sheet, x=1cm, y=1cm]
  \node[cg=sheetRed] (root) at (0,0) {/};
  \node[why] at (2.6,0) {all of it \textbf{read-only}};
  \node[dir] at (5.05,0) {ProtectSystem=strict};
  \node[cg=sheetGrey] (h) at (0.3,-0.27) {/home /root /run/user};
  \node[why] at (2.6,-0.27) {empty};
  \node[dir] at (5.05,-0.27) {ProtectHome=yes};
  \node[cg=sheetBlue] (t) at (0.3,-0.54) {/tmp /var/tmp};
  \node[why] at (2.6,-0.54) {its own};
  \node[dir] at (5.05,-0.54) {PrivateTmp=yes};
  \node[cg=sheetBlue] (d) at (0.3,-0.81) {/dev};
  \node[why] at (2.6,-0.81) {null, zero, random — no disks};
  \node[dir] at (5.05,-0.81) {PrivateDevices=yes};
  \node[cg=sheetGreen!60!black] (s) at (0.3,-1.08) {/var/lib/app};
  \node[why] at (2.6,-1.08) {\textbf{writable}, owned by it};
  \node[dir] at (5.05,-1.08) {StateDirectory=app};
  \node[cg=sheetBlue] (n) at (0,-1.35) {network};
  \node[why] at (2.6,-1.35) {\ct{lo} only};
  \node[dir] at (5.05,-1.35) {PrivateNetwork=yes};
  \foreach \c in {h,t,d,s} \draw[sheetGrey] (0.1,-0.1) |- (\c.west);
\end{tikzpicture}\par
Each row is a namespace trick: the service gets its own mount (and network) namespace.
\ct{NoNewPrivileges=}: \ct{execve} ignores setuid ·
\ct{SystemCallFilter=@system-service}: other syscalls SIGSYS.
\textbf{\link{\ct{DynamicUser=}}{https://man7.org/linux/man-pages/man5/systemd.exec.5.html}} (232)
leases a UID from \textbf{61184–65519} (\ct{0xEF00}–\ct{0xFFEF}: above the users distributions hand
out, below 65534 \ct{nobody}) while the service runs, and implies most rows above. A recycled UID
must not inherit files, so state lives in \ct{/var/lib/private/app}, behind a root-only \ct{0700}
directory. \ct{systemd-analyze security}: 0–10, high = exposed.
\textbf{Secrets}: \ct{Environment=} is ``exposed to unprivileged clients via D-Bus'';
\ct{LoadCredential=} (247): read-only, unswappable, only the unit's user; \ct{…Encrypted=} (250):
sealed to TPM2 / a host key.

\hd{journald — logs as fields, some unforgeable}
stdout / stderr become entries of fields (\ct{MESSAGE}, \ct{\_PID}, \ct{\_SYSTEMD\_UNIT} …). A
leading underscore marks a
\link{trusted field}{https://man7.org/linux/man-pages/man7/systemd.journal-fields.7.html}: journald
adds it, \textbf{client code cannot alter it} — a process can lie in \ct{MESSAGE}, not in \ct{\_PID}.
\textbf{\link{Forward Secure Sealing}{https://man7.org/linux/man-pages/man1/journalctl.1.html}}:
\ct{journalctl \dd setup-keys} leaves a sealing key on the host that moves on, one way, every
15\,min, and a verification key kept elsewhere — an intruder can rewrite only the current window
unnoticed by \ct{\dd verify}: a 2013 cryptography paper (``seekable sequential key generators'')
shipped in the log daemon. Persistent storage: default since 259; disk 10\,\% used /
15\,\% free ($\le$4\,G); 10\,000 messages / 30\,s per service.

\hd{Releases}
\noindent\begin{tikzpicture}[sheet, x=1cm, y=1cm, ev/.style={font=\tiny, align=center, inner sep=0.5pt}]
  \draw[thick, sheetGrey] (0,0) -- (7.9,0);
  \foreach \x/\v/\d/\t/\p in {
      0.55/247/2020-11/{credentials}/n,
      1.3/250/2021-12/{encrypted\\credentials}/s,
      2.1/254/2023-07/{soft-reboot,\\confext}/n,
      2.95/256/2024-06/{run0; refuses\\cgroup v1}/s,
      3.95/258/2025-09/{cgroup v1, telinit,\\runlevel removed}/n,
      4.95/259/2025-12/{journal\\persistent}/s,
      6.0/260/2026-03/{SysV scripts,\\rc-local removed}/n,
      7.3/262/2026-09/{latest}/s} {
    \fill[sheetOrange] (\x,0) circle (0.6mm);
    \if\p n \node[ev, anchor=south] at (\x,0.07) {\textbf{\v} · \d\\\t};
    \else \node[ev, anchor=north] at (\x,-0.07) {\textbf{\v} · \d\\\t}; \fi }
\end{tikzpicture}

\begin{htmlonly}
\section{Further reading}
\begin{itemize}
  \item \href{https://0pointer.de/blog/projects/systemd.html}{Rethinking PID 1} (2010) — the essay
        that started it: why socket activation makes ordering unnecessary, and what it took from
        inetd and launchd
  \item \href{https://0pointer.net/blog/dynamic-users-with-systemd.html}{Dynamic Users with systemd}
        (2017) — why the UID range is 61184–65519 and how \texttt{/var/lib/private} survives UID recycling
  \item \href{https://eprint.iacr.org/2013/397}{Practical Secure Logging: Seekable Sequential Key Generators}
        (2013) — the cryptography paper behind \ct{journalctl \dd setup-keys}, which became part of journald
  \item \href{https://systemd.io/CGROUP_DELEGATION/}{Control Group APIs and Delegation} — the
        single-writer rule, and what a delegated service must do with its own subtree
  \item \href{https://0pointer.de/blog/projects/the-biggest-myths.html}{The Biggest Myths} (2013) —
        the early arguments over systemd, answered point by point
\end{itemize}
\end{htmlonly}

\end{multicols}

\noindent{\raggedright\footnotesize\color{sheetGrey}\textit{Related:} \texttt{linux-boot-process}
(targets, PID~1, rescue) · \texttt{linux-memory-oom} (\ct{Memory*=}, systemd-oomd) ·
\texttt{namespaces-cgroups-containers} (cgroup v2) · \texttt{linux-security-permissions} (caps,
seccomp, run0) · \texttt{linux-performance-observability} · \texttt{os-fundamentals} ·
\texttt{xnu-launchd-xpc}\par}

\end{document}
